\uuid{FEv8}
\exo7id{1582}
\titre{exo7 1582}
\auteur{ortiz}
\organisation{exo7}
\datecreate{1999-04-01}
\isIndication{false}
\isCorrection{false}
\chapitre{Réduction d'endomorphisme, polynôme annulateur}
\sousChapitre{Polynôme annulateur}
\cpge{mpsi-mp2i.polynomes-fractions-rationnelles.a, pcsi.polynomes.a, ptsi.polynomes.a, mpsi-mp2i.applications-lineaires.b, pcsi.applications-lineaires.b, ptsi.applications-lineaires.b, mp-mpi.reduction-endomorphismes.g, pc.complements-espaces-vectoriels.d, psi.complements-espaces-vectoriels.d, mp-mpi.reduction-endomorphismes.i, pc.reduction.d, psi.reduction-endomorphismes.d}
\module{Algèbre}
\niveau{L2}
\difficulte{}

\contenu{
\question{
Que peut-on dire d'un endomorphisme d'un
$K$-espace vectoriel de dimension finie annul\'e
par les polyn\^omes $P=1-X^3$ et $Q=X^2-2X+1$?
}
}
